% Auxiliary to OpenAI, A two-dimensional area law from a global spectral gap,
% Section 11, geometry:initial-stars, lines 352--370, and
% prop:two-families, lines 308--323;
% revision adc7f1241b42e322a6451854ab7e4b4c146bf78a.

\section{Connected components after radial deletion}
\label{sec:area_law_fan_cut_components}

The cut fan pieces of Definition~\ref{def:area_law_fan_cut_piece}
combine along each equal-color run to form a connected component of the cut
square, as used in the initial-star construction of
\cite[Section~11]{OpenAI2026AreaLaw}.

\begin{theorem}[Connected components of a cut fan square]
    \label{thm:area_law_fan_cut_component_correspondence}
    \lean{TNLean.PEPS.AreaLaw.Geometry.cellFanRun_connectedComponents_equiv}
    \leanok
    \uses{def:area_law_fan_cut_piece,def:area_law_fan_runs}
    Let $Q$, its midpoint-subdivided fan, its two-coloring, and the union of
    separating radial segments $L$ be as in Definition~\ref{def:area_law_fan_cut_piece}.
    Keep an adjacency edge between consecutive triangles when their colors
    agree, and let $\mathcal R$ be the set of connected components (runs) of
    this finite graph. Set $\Omega=B_\infty(c,r)\setminus L$, and for each
    $R\in\mathcal R$ define
    \begin{align}
        T_R & =\bigcup_{i\in R}P_i, \\
        U_R & =T_R\cap\Omega.
    \end{align}
    Write $\mathcal C(\Omega)$ for the connected components of $\Omega$
    and $C_\Omega(x)$ for the component containing $x\in\Omega$.
    Then there is a unique bijection
    $\Phi:\mathcal C(\Omega)\longrightarrow\mathcal R$ such that, for
    every $x\in\Omega$ and $R\in\mathcal R$,
    \begin{align}
        \Phi(C_\Omega(x))=R & \quad\Longleftrightarrow\quad x\in T_R.
        \label{eq:area_law_fan_cut_component_membership}
    \end{align}
    Equivalently, the map $R\mapsto U_R$ is a bijection from $\mathcal R$ onto
    $\mathcal C(\Omega)$.
\end{theorem}
\begin{proof}
    \leanok
    \uses{def:area_law_fan_runs,
        thm:area_law_sample_rows,
        thm:area_law_cell_fan_closed_ball,
        thm:area_law_fan_run_cover,
        thm:area_law_fan_cut_runs_disjoint,
        thm:area_law_fan_cut_run_connected}
    By Theorem~\ref{thm:area_law_cell_fan_closed_ball},
    $Q=\overline B_\infty(c,r)$, so $\Omega\subset Q$.
    By Theorem~\ref{thm:area_law_fan_run_cover}, the sets $T_R$ cover
    $Q$, so the sets $U_R$ cover $\Omega$.
    By Theorem~\ref{thm:area_law_fan_cut_runs_disjoint}, the sets $U_R$ are
    pairwise disjoint. Since $U_R=(T_R\setminus L)\cap B_\infty(c,r)$,
    Theorem~\ref{thm:area_law_fan_cut_run_connected} shows that every
    $U_R$ is nonempty and connected in $\Omega$.

    Each triangle $P_i$ is compact and closed, so each finite union $T_R$ is
    closed in $\mathbb R^2$ and $U_R=T_R\cap\Omega$ is closed in $\Omega$.
    Since $\mathcal R$ is finite and the sets $U_R$ partition $\Omega$, the
    complement $\Omega\setminus U_R=\bigcup_{T\ne R}U_T$ is also closed in
    $\Omega$, so each $U_R$ is both open and closed in $\Omega$.

    Every nonempty connected subset of $\Omega$ that is both open and closed in
    $\Omega$ is a connected component. Thus $\{U_R:R\in\mathcal R\}$ is the set
    of connected components of $\Omega$, and for $x\in\Omega$ we have
    $x\in U_R$ if and only if $x\in T_R$, proving
    (\ref{eq:area_law_fan_cut_component_membership}).
\end{proof}
